\documentclass{amsart}
% For dealing with references we use the comment environment
\usepackage{verbatim}
\newenvironment{reference}{\comment}{\endcomment}
%\newenvironment{reference}{}{}
\newenvironment{slogan}{\comment}{\endcomment}
\newenvironment{history}{\comment}{\endcomment}

% For commutative diagrams we use Xy-pic
\usepackage[all]{xy}

% We use 2cell for 2-commutative diagrams.
\xyoption{2cell}
\UseAllTwocells

% We use multicol for the list of chapters between chapters
\usepackage{multicol}

% This is generally recommended for better output
\usepackage{lmodern}
\usepackage[T1]{fontenc}

% For cross-file-references

% Package for hypertext links:
\usepackage{hyperref}

% For any local file, say "hello.tex" you want to link to please

% Theorem environments.
%
\theoremstyle{plain}
\newtheorem{theorem}[subsection]{Theorem}
\newtheorem{proposition}[subsection]{Proposition}
\newtheorem{lemma}[subsection]{Lemma}

\theoremstyle{definition}
\newtheorem{definition}[subsection]{Definition}
\newtheorem{example}[subsection]{Example}
\newtheorem{exercise}[subsection]{Exercise}
\newtheorem{situation}[subsection]{Situation}

\theoremstyle{remark}
\newtheorem{remark}[subsection]{Remark}
\newtheorem{remarks}[subsection]{Remarks}

\numberwithin{equation}{subsection}

% Macros
%
\def\lim{\mathop{\mathrm{lim}}\nolimits}
\def\colim{\mathop{\mathrm{colim}}\nolimits}
\def\Spec{\mathop{\mathrm{Spec}}}
\def\Hom{\mathop{\mathrm{Hom}}\nolimits}
\def\Ext{\mathop{\mathrm{Ext}}\nolimits}
\def\SheafHom{\mathop{\mathcal{H}\!\mathit{om}}\nolimits}
\def\SheafExt{\mathop{\mathcal{E}\!\mathit{xt}}\nolimits}
\def\Sch{\mathit{Sch}}
\def\Mor{\mathop{\mathrm{Mor}}\nolimits}
\def\Ob{\mathop{\mathrm{Ob}}\nolimits}
\def\Sh{\mathop{\mathit{Sh}}\nolimits}
\def\NL{\mathop{N\!L}\nolimits}
\def\CH{\mathop{\mathrm{CH}}\nolimits}
\def\proetale{{pro\text{-}\acute{e}tale}}
\def\etale{{\acute{e}tale}}
\def\QCoh{\mathit{QCoh}}
\def\Ker{\mathop{\mathrm{Ker}}}
\def\Im{\mathop{\mathrm{Im}}}
\def\Coker{\mathop{\mathrm{Coker}}}
\def\Coim{\mathop{\mathrm{Coim}}}

% Boxtimes
%
\DeclareMathSymbol{\boxtimes}{\mathbin}{AMSa}{"02}

%
% Macros for moduli stacks/spaces
%
\def\QCohstack{\mathcal{QC}\!\mathit{oh}}
\def\Cohstack{\mathcal{C}\!\mathit{oh}}
\def\Spacesstack{\mathcal{S}\!\mathit{paces}}
\def\Quotfunctor{\mathrm{Quot}}
\def\Hilbfunctor{\mathrm{Hilb}}
\def\Curvesstack{\mathcal{C}\!\mathit{urves}}
\def\Polarizedstack{\mathcal{P}\!\mathit{olarized}}
\def\Complexesstack{\mathcal{C}\!\mathit{omplexes}}
% \Pic is the operator that assigns to X its picard group, usage \Pic(X)
% \Picardstack_{X/B} denotes the Picard stack of X over B
% \Picardfunctor_{X/B} denotes the Picard functor of X over B
\def\Pic{\mathop{\mathrm{Pic}}\nolimits}
\def\Picardstack{\mathcal{P}\!\mathit{ic}}
\def\Picardfunctor{\mathrm{Pic}}
\def\Deformationcategory{\mathcal{D}\!\mathit{ef}}

\setlength{\emergencystretch}{3em}
\begin{document}
\title[Stacks Project, Tag 0G3L]{413 --- Controlled Ext factorizations for powers of an ideal}
\maketitle
\noindent Copyright (C) 2005--2025 Johan de Jong. Stacks Project authors. Source: \href{https://stacks.math.columbia.edu/tag/0G3L}{Tag 0G3L}.

\noindent Editorial difficulty note (not author text). Difficulty H2: The author chooses compatible presentations for a module and its ideal-power submodules and controls the syzygy comparison by Artin-Rees. A diagram proves the first-Ext factorization, and dimension shifting inductively supplies a uniform loss of ideal powers in every higher degree..

\noindent Editorial provenance note (not author text). History: excerpt from revision \texttt{a0444\allowbreak 6e57e\allowbreak c1fbc\allowbreak 252a8\allowbreak 71afc\allowbreak ec775\allowbreak 2fb28\allowbreak 07b14}, more-algebra.tex, lines 29492--29568. Reference presentation was adapted; original-author context and core support blocks were retained or added. The mathematical wording is unchanged. Licensed under GNU FDL 1.2 or later; no invariant sections or cover texts. See ../licenses/COPYING. Context blocks reproduce author definitions and supporting results; they are not additional counted problems.

\begin{lemma}
\label{lemma-ext-factors}
Let $A$ be a Noetherian ring. Let $I \subset A$ be an ideal.
Let $M$, $N$ be $A$-modules with $M$ finite. For each $p > 0$ there exists a
$c \geq 0$ such that for $n \geq c$ the map
$\Ext_A^p(M, N) \to \Ext_A^p(I^nM, N)$
factors through $\Ext^p_A(I^nM, I^{n - c}N) \to \Ext_A^p(I^nM, N)$.
\end{lemma}

\begin{proof}
For $p = 0$, if $\varphi : M \to N$ is an $A$-linear map, then
$\varphi(\sum f_i m_i) = \sum f_i \varphi(m_i)$ for $f_i \in A$
and $m_i \in M$. Hence $\varphi$ induces a map $I^nM \to I^nN$
for all $n$ and the result is true with $c = 0$.

\medskip\noindent
Choose a short exact sequence $0 \to K \to A^{\oplus t} \to M \to 0$.
For each $n$ we pick a short exact sequence
$0 \to L_n \to A^{\oplus s_n} \to I^nM \to 0$.
It is clear that we can construct a map of short exact sequences
$$
\xymatrix{
0 \ar[r] &
L_n \ar[r] \ar[d] &
A^{\oplus s_n} \ar[r] \ar[d] &
I^nM \ar[r] \ar[d] & 0 \\
0 \ar[r] &
K \ar[r] &
A^{\oplus t} \ar[r] &
M \ar[r] & 0
}
$$
such that $A^{\oplus s_n} \to A^{\oplus t}$ has image in $(I^n)^{\oplus t}$.
By Artin-Rees (Algebra, Lemma \href{https://stacks.math.columbia.edu/tag/00IN}{lemma Artin Rees (Tag 00IN)}) there exists
a $c \geq 0$ such that $L_n \to K$ factors through $I^{n - c}K$ if $n \geq c$.

\medskip\noindent
For $p = 1$ our choices above induce a solid commutative diagram
$$
\xymatrix{
\Hom_A(A^{\oplus s_n}, N) \ar[r] &
\Hom_A(L_n, N) \ar[r] &
\Ext_A^1(I^nM, N) \ar[r] & 0 \\
\Hom_A((I^n)^{\oplus t}, I^{n - c}N) \ar[r] \ar[u] &
\Hom_A(K \cap (I^n)^{\oplus t}, I^{n - c}N) \ar[r] \ar[u] &
\Ext_A^1(I^nM, I^{n - c}N) \ar[u] \\
\Hom_A(A^{\oplus t}, N) \ar[r] \ar[u] &
\Hom_A(K, N) \ar[r] \ar[u] &
\Ext_A^1(M, N) \ar@{..>}[u] \ar[r] & 0
}
$$
whose horizontal arrows are exact. The lower middle vertical arrow
arises because $K \cap (I^n)^{\oplus t} \subset I^{n - c}K$
and hence any $A$-linear map $K \to N$ induces an $A$-linear map
$(I^n)^{\oplus t} \to I^{n - c}N$ by the argument of the first paragraph.
Thus we obtain the dotted arrow as desired.

\medskip\noindent
For $p > 1$ we obtain a commutative diagram
$$
\xymatrix{
\Ext^{p - 1}_A(I^{n - c}K, N) \ar[r] &
\Ext^{p - 1}_A(L_n, N) \ar[r] &
\Ext_A^p(I^nM, N) \\
\Ext^{p - 1}_A(K, N) \ar[rr] \ar[u] & &
\Ext_A^p(M, N) \ar[u]
}
$$
whose bottom horizontal arrow is an isomorphism. By induction
on $p$ the left vertical map factors through
$\Ext^{p - 1}_A(I^{n - c}K, I^{n - c - c'}N)$ for some $c' \geq 0$
and all $n \geq c + c'$. Using the composition
$\Ext^{p - 1}_A(I^{n - c}K, I^{n - c - c'}N) \to
\Ext^{p - 1}_A(L_n, I^{n - c - c'}N) \to \Ext^p_A(I^nM, I^{n - c - c'}N)$
we obtain the desired factorization (for $n \geq c + c'$ and with $c$
replaced by $c + c'$).
\end{proof}
\end{document}
